% main.tex -- manuscript. SOURCE OF TRUTH from 2026-10-15 (docs/manuscript/DRAFT.md is a frozen snapshot).
% Build: ./build.sh (tectonic). Figures are read from ../docs/manuscript/figures/ (see macros.tex).
% LAYER: SKELETON (2026-10-26). Built on docs/manuscript/OUTLINE.md draft 3. Each paragraph holds its first sentence
% (the claim) only; the comment under it lists what the paragraph will carry and where each number comes from.
% Prose comes in the next layer, after the author has reacted to this skeleton. The previous prose draft (120-day
% story) is in git history (commit 9446e40 and earlier).
\documentclass[largeformat]{interact} % Taylor & Francis template (CMBBE: Imaging & Visualization)
\input{macros}

\begin{document}
\linenumbers

\articletype{ARTICLE}

\title{\papertitle}

\author{
\name{\todo{PENDING: first author's full name}\textsuperscript{a}\thanks{CONTACT \todo{PENDING: first author. Email: \ldots; ORCID: \ldots}}
  and Jessica Zhang\textsuperscript{b}}
\affil{\textsuperscript{a}Department of Biomedical Engineering, Carnegie Mellon University, Pittsburgh, PA, USA;
\textsuperscript{b}Department of Mechanical Engineering, Carnegie Mellon University, Pittsburgh, PA, USA}
}

\maketitle

% Unstructured abstract, 100--150 words (journal instructions).
% PROPOSED 2026-10-26 (149 words), to replace the abstract locked on 2026-10-24 (docs/manuscript/LOCKED.md), which
% predates the external tests and the system results. Not locked until the author approves it.
\begin{abstract}
\flag{[PROPOSED, awaiting approval]} Models of post-treatment glioma MRI face known challenges: most datasets lack
treatment dates, centres document progression differently, cohorts are small and the hardest situations have few
cases. We built PRISM, a modular system that keeps image and treatment-timeline evidence in separate pathways,
measures the tumour by automatic segmentation, attaches a confidence to every scan and writes a report with a local
language model. On MU-Glioma-Post, the timeline pathway detected the scan at which progression was called with an
AUC of 0.82 and the image pathway with 0.71. Fitted once and locked, the models reached 0.77 on LUMIERE, and
single-scan images reached 0.80 on UCSF-ALPTDG, about as high as models trained there. An audit showed what each
score rested on: the treatment timeline at MU, enhancing tumour volume at UCSF. Conformal sets kept their coverage at
all three centres, and every evaluation report passed a deterministic check against its record.
\end{abstract}

% No more than 3 keywords (journal instructions).
\begin{keywords}
glioma surveillance; external validation; clinical decision support
\end{keywords}

% =====================================================================================================
\section{Introduction}\label{sec:1}

After surgery and chemoradiotherapy, patients with glioma are followed with MRI every few months, and at each scan
the team must decide whether the tumour has progressed.
% Why this is hard: enhancement can grow after RT without tumour growth (pseudoprogression, most common in the first
% weeks after chemoradiotherapy, more frequent with MGMT methylation) \citep{brandes2008,taal2008}; progression in
% routine records is documented from imaging, symptoms and the scan schedule. No RANO.

Models for this task meet five challenges that come from the data rather than from any one method.
% C1 most datasets have images but no treatment dates (LUMIERE, UCSF);
% C2 centres document progression differently, so a score can rest on different evidence;
% C3 cohorts are small and a single split can move a score more than a model choice \citep{an2021};
% C4 the hardest situations (early after RT, growing enhancement) have few cases;
% C5 volumes are noisy, the clinician must see what was measured, and patient data cannot leave the hospital.
% External tests are uncommon in this field (\todo{cite Shi 2025, Neuro-Oncology: 3 of 20}; \citealp{yu2022}).
% Neutral: properties of the data, not a judgement of earlier work.

We asked whether a system could give the clinician, at each post-treatment scan, the image evidence and the
treatment-timeline evidence separately, carry to other centres without refitting, and say how far each can be trusted
there.

PRISM is a modular system built for this, and we tested it on one development cohort and two external cohorts with
models fitted once and plans written in advance.
% One sentence per module (Fig 1) and the headline findings in words: the pathways transport; the audit shows what
% each score rests on; the confidence keeps its coverage; the reports are faithful to the record.

% =====================================================================================================
\section{Methods}\label{sec:2}

\subsection{PRISM}\label{sec:2.1}

PRISM handles each scan in five modules, each answering one of the challenges above (Figure~\ref{fig:1}).
% Module list; the challenge -> module pairing is Table 2 in the outline (C1--C5).

The timeline pathway scores a scan from clinical data and the patient's position in treatment.
% Age, sex, IDH, MGMT; days since diagnosis and surgery; RT start and end; chemotherapy; previous scans. Ridge,
% patient-grouped. On cohorts without treatment dates, the transportable version (days since surgery, early post-RT
% flag). Sources: DETECTION_TASK.md (B1, Tbase).

The image pathway scores a scan from simple measurements of the segmented tumour.
% Volumes (enhancing, necrosis, FLAIR hyperintensity, cavity), shape, relative intensity; log ratios to the previous
% scan and to the smallest earlier volume. No timing inputs. Sources: DETECTION_TASK.md (I_scan, I_pure).

A deep segmentation network supplies these measurements and the images shown in the report.
% BraTS21 SwinUNETR (MU, LUMIERE: the cohorts' own masks or resegmentation; UCSF: provided masks) -- check per cohort
% in UCSF_IMAGING.md and RESEGMENTATION.md. The segmentation figure is drawn by code; the language model never sees
% an image. Source: SEG_FIGURE.md.

Each scan's confidence has three parts: a conformal set, a pathway-agreement statement and a reliability card.
% Patient-level split conformal, alpha 0.2 (and 0.1), calibrated on MU; agreement = both scores on the same side of
% their cohort medians; card = measured AUCs by situation (phase x enhancement), tier tested / few cases / not tested.
% Sources: SYSTEM_PLAN.md, RELIABILITY_CARD.md.

A local language model writes the report from the scan's record, and a deterministic checker tests every sentence
against that record.
% Two sentences only in the main text (decision 6). Qwen3.8-27B, offline, temperature 0; checker: numbers, flags,
% conformal and agreement phrases, forbidden wording (diagnosis, advice, probability words). Detail: supplement.

\subsection{Cohorts and the task}\label{sec:2.2}

PRISM was developed on MU-Glioma-Post and tested, locked, on LUMIERE and UCSF-ALPTDG (Table~\ref{tab:1}).
% \citep{mahmoud2025,suter2022}, \todo{cite UCSF-ALPTDG}. What each cohort has: images, treatment dates, how
% progression is documented. UCSF: one scan pair per patient, so the evaluation is restricted to scan 2
% (UCSF_PLAN.md). UCSF obtained under a data use agreement.

The task is detection at the scan: is this the scan at which progression is called?
% Positives and negatives as in DETECTION_TASK.md and UCSF_PLAN.md; grade 4 primary; the intended use on MU (scans
% from the start of RT). 120-day risk as the secondary task, one sentence.

\subsection{Evaluation}\label{sec:2.3}

On MU, every model was evaluated by repeated patient-grouped cross-validation, because a single split moved scores
more than any model choice (Supplementary Figure~\todo{S-split}).
% 5 x 5, patient-cluster bootstrap 1000 draws. Supporting evidence in one or two sentences: MU first split vs 500
% splits (0.80 vs 0.72); UCSF's own split vs 100 random splits. The single-scan radiomics setting goes to the
% supplement (CHRISTODOULOU_REPLICATION.md).

For the external tests, each model was fitted once on MU and applied once, under plans committed before any external
label was read.
% UCSF_PLAN.md (0e7c9fe), SYSTEM_PLAN.md (e7cba80), REPORT_V3_1_PLAN.md (4f7e7b9). Questions fixed in advance and
% their answers: Table~\ref{tab:2}.

Analyses that were not planned are labelled exploratory wherever they appear.
% The image-model ladder; the reduced confidence design; the report v3.1 run (a second evaluation after a seen
% failure); within-phase and within-enhancement comparisons are diagnostics (decision 8).

% =====================================================================================================
\section{Results}\label{sec:3}

\subsection{Each pathway carries evidence on MU-Glioma-Post}\label{sec:3.1}

In the intended use on MU-Glioma-Post, the timeline pathway detected the scan at which progression was called with an
AUC of 0.82, and the image pathway with 0.71.
% Combined 0.82; clinical data alone at chance (0.53). Interval on the headline numbers only. Source:
% DETECTION_TASK.md, exploratory images + clinical section.

\subsection{The pathways hold at two external centres}\label{sec:3.2}

Fitted once on MU and applied unchanged, the combined model reached 0.77 on LUMIERE, and single-scan image
measurements reached 0.80 on UCSF-ALPTDG (Figure~\ref{fig:2}).
% LUMIERE: images alone 0.66, timeline alone 0.69, both 0.77 (complementary). UCSF: models trained on UCSF reach 0.83
% for single-scan images, so the locked model is about as high (claim transport, not better performance; decision 10).
% Clinical alone at chance in all cohorts. Sources: DETECTION_TASK.md, UCSF_RESULTS.md.

\subsection{What each score rests on depends on the cohort}\label{sec:3.3}

The audit shows that on MU the scores rest mainly on where the patient is in treatment, and on UCSF on how much the
tumour enhances (Figure~\ref{fig:3}).
% MU: the timeline's gain lies between phases of care (+0.11 overall, -0.02 within phase; diagnostic).
% UCSF: enhancing volume alone 0.88; 27 of 29 progression scans were rated enhancement increased.

Each step of the image model's design raised its AUC on MU, and the audit shows what each step added.
% EXPLORATORY ladder: single scan 0.66 -> + change 0.71 -> + schedule context 0.81; the last step encodes the
% timeline. One sentence on the deep networks: as segmenters they supply PRISM's measurements; their direct scores
% track enhancing volume (r about 0.8) -- framed as segmenters, not failed classifiers (STYLE 2026-10-26).

\subsection{The confidence keeps its coverage at all three centres}\label{sec:3.4}

At the 80\% level, the conformal sets of the combined score covered 0.82 of MU scans, 0.87 of LUMIERE scans and 0.78
of UCSF scans, at a nominal level of 0.80 and a threshold of 0.78 fixed in advance.
% Figure~\ref{fig:4}A. Per-pathway coverage in the supplement (images on LUMIERE and timeline on UCSF below target).

When the two pathways disagreed on MU, the combined score was less accurate, and setting those scans aside raised
the AUC of the rest more than removing the same number at random.
% S2 +0.31, S3 +0.08 (Figure~\ref{fig:4}B). Disagreement does not pick out errors on UCSF. The reduced design
% (disagreement AND conformal set with both outcomes) is EXPLORATORY: flags 12--27\%, wrong about half the time vs a
% quarter otherwise, in all three cohorts (SYSTEM_RESULTS.md).

For most scans the evidence behind any score is thin, and the reliability card says so.
% 35 of 249 MU scans in a situation tested on >= 10 progression and >= 10 non-progression scans; every UCSF situation
% has few cases; the card states evidence, not accuracy (Figure~\ref{fig:4}C). Source: RELIABILITY_CARD.md.

\subsection{The reports are faithful to the record}\label{sec:3.5}

Every evaluation report written by the local model passed the deterministic check against its record: 114 of 114 on
MU, 162 of 162 on LUMIERE and 130 of 130 on UCSF (Figure~\ref{fig:5}).
% Disclose: second evaluation after a failure seen in the first (the LUMIERE rater flag; first evaluation 0 of 178),
% with the checker unchanged (REPORT_V3_1.md). Tier stated wherever evidence is thin; reports followed a changed
% field in 28 of 30. Models smaller than 27B left required content out (ladder in the supplement).

% =====================================================================================================
\section{Discussion}\label{sec:4}

PRISM gives the clinician separate image and treatment-timeline evidence at each scan, carries to two external
centres without refitting, and states how far its scores can be trusted there.
% The answer to the question, first sentence.

Because its modules are separate, PRISM runs where data are incomplete.
% The image pathway works without treatment dates (LUMIERE, UCSF); the timeline is added without a new model; a new
% cohort adds rows to the card.

Consider a scan eight weeks after radiotherapy in which the enhancing tumour has grown.
% Scenario: the pathways disagree, the conformal set holds both outcomes, the card says this situation has few
% cases. PRISM states this; it does not decide.

Each of PRISM's components has been described before, but we found no system that combines them.
% Concise, neutral (SYSTEM_LANDSCAPE.md): segmentation and volumetrics; clinical + volumetric fusion (Jabal 2026);
% conformal prediction in imaging; model cards. Avoid "first". Literature comparison: Christodoulou 2026 (label at
% the patient level, one held-out set; one neutral sentence) \citep{christodoulou2026}, Tikhonov 2025, Chen 2026
% \citep{chen2026}.

This study has limits.
% One development centre; documented, not expert-read, progression; small subgroups; the ladder and the reduced
% confidence design are exploratory; the report checker is lexical; no clinician has read the reports.

The next step is to pool the rare situations across cohorts and train on them.
% Long tail: early post-RT progression, growing enhancement without progression; Burdenko and UCSD-PTGBM; a clinician
% reader study.

\subsection*{Conclusion}

PRISM shows that a surveillance system can keep its kinds of evidence apart, carry them to other centres and say
where they have and have not been shown to work.

% =====================================================================================================
\section*{Acknowledgements}
\todo{PENDING AUTHOR DECISION}.

\section*{Disclosure statement}
\todo{PENDING AUTHOR DECISION: competing interests, or ``The authors report there are no competing interests to
declare.''} If any author serves on the editorial or advisory board of this journal, state it here.

\section*{Declaration of generative AI use}
\todo{DRAFT, PENDING AUTHOR CONFIRMATION:} Generative AI is part of the system studied. PRISM's report generator uses
a local language model (Qwen3.8-27B; six other open-weight models in a comparison); all ran offline on the authors'
computing cluster and saw only derived system records. An AI assistant (Claude, Anthropic) was used to write and run
analysis code and to draft and edit the manuscript under the authors' direction; it was given no patient-level
data through any hosted service. The authors reviewed all code, results and text and take full responsibility for
the content of this article.

\section*{Funding}
\todo{PENDING AUTHOR DECISION: ``This work was supported by the [Funding Agency] under Grant [number].'' If there was
no funding: ``No funding was obtained for this work.''}

\section*{Ethics statement}
\todo{PENDING AUTHOR DECISION: secondary analysis of de-identified datasets (MU-Glioma-Post and LUMIERE, public;
UCSF-ALPTDG under a data use agreement); state the institutional determination or waiver and cite the source
datasets' approvals.} No registered protocol. Analysis plans were committed to the repository before each external
test. Patient and public involvement: none.

\section*{Data availability statement}
\todo{PENDING AUTHOR DECISION: MU-Glioma-Post (TCIA, doi:10.7937/\allowbreak 7K9K-3C83) and LUMIERE (figshare;
\citealp{suter2022}) are public; UCSF-ALPTDG is available from its provider under a data use agreement; state which
derived tables, fold assignments and locked model files are released, with a DOI.} Code: \todo{PENDING: public
repository link and release tag.}

\section*{Supplemental material}
Supplemental data for this article are provided as a separate file.

\section*{Notes on contributors}
\todo{PENDING: a short biographical note for each author (no more than 200 words each).}

\bibliography{refs}

% =====================================================================================================
\clearpage
\section*{Tables}

\begin{table}[!htbp]
\centering
\caption[The three cohorts]{\textbf{The three cohorts.} \todo{SKELETON: per cohort, what it has (images, treatment
dates, how progression is documented), patients, scans in the detection task, positives and negatives; grade 4
primary. Sources: TABLE1\_COHORT.md (MU, LUMIERE), UCSF\_DATA.md.}}
\label{tab:1}
\end{table}

\begin{table}[!htbp]
\centering
\caption[Questions fixed in advance and their answers]{\textbf{Questions fixed in advance and their answers.}
\todo{SKELETON: MU detection (DETECTION\_TASK.md), LUMIERE, UCSF Q1--Q4 (UCSF\_PLAN.md), system S1--S7
(SYSTEM\_PLAN.md, REPORT\_V3\_1\_PLAN.md); each with its rule, result and met / not met.}}
\label{tab:2}
\end{table}

% =====================================================================================================
\clearpage
\section*{Figure legends}

\begin{figure}[!htbp]
\centering
\figslot{fig1_system.pdf}{0.40\textheight}{PRISM is built of modules, each answering a known challenge: scan and
record; segmentation and measurements; the timeline pathway; the image pathway; the combined score; the confidence
(conformal set, agreement, card); the local report. One worked example scan. No RANO.}
\caption{\textbf{PRISM is built of modules, each answering a known challenge.} \todo{legend}}
\label{fig:1}
\end{figure}

\begin{figure}[!htbp]
\centering
\figslot{fig2_transport.pdf}{0.40\textheight}{AUC by pathway (clinical alone, timeline, images, combined) and cohort
(MU internal; LUMIERE and UCSF locked; UCSF-trained as a reference).}
\caption{\textbf{The pathways hold at two external centres without refitting.} \todo{legend}}
\label{fig:2}
\end{figure}

\begin{figure}[!htbp]
\centering
\figslot{fig3_evidence.pdf}{0.40\textheight}{(A) MU: the timeline's gain lies between phases of care. (B) UCSF:
enhancing volume alone matches the image model. (C) The image-model ladder (exploratory).}
\caption{\textbf{What each score rests on depends on the cohort.} \todo{legend}}
\label{fig:3}
\end{figure}

\begin{figure}[!htbp]
\centering
\figslot{fig4_confidence.pdf}{0.40\textheight}{(A) Conformal coverage at the 80\% level by cohort and pathway. (B) AUC
with and without the disagreement warning. (C) The reliability card: situations by cohort with their evidence tier.}
\caption{\textbf{The confidence keeps its coverage and points to doubtful scans.} \todo{legend}}
\label{fig:4}
\end{figure}

\begin{figure}[!htbp]
\centering
\figslot{fig5_report.pdf}{0.40\textheight}{One rendered report (segmentation figure, pathways, combined score,
confidence) and faithfulness by cohort.}
\caption{\textbf{The report states the evidence, faithful to the record.} \todo{legend}}
\label{fig:5}
\end{figure}

\end{document}
